\section{Conclusion and Discussion}
We have introduced a new family of video foundation models called {\modelname}, which achieves the state-of-the-art performance across various video and audio tasks. In {\modelname}, we combine masked video modeling, video-audio-text contrastive learning, and next token prediction into a unified framework. Additionally, we create a new video-text dataset that incorporates video-audio-speech fused captions as descriptions. The dataset contains temporally segmented clips with high semantic coherence.
These designs in {\modelname} contribute to enhancing video understanding in both perception and reasoning tasks. Notably, {\modelname} excels in video-related dialogue and long video understanding, demonstrating its effectiveness in capturing high-level semantics.

\paragraph{Limitations and Discussions.} Despite its achievements, {\modelname} does not introduce specific novel architectural design. Instead, it leverages the existing learning techniques for scaling video foundation models while focusing on improving data processing to enhance its spatiotemporal perception, semantic alignment, and basic knowledge embedding. Similarly to previous studies~\cite{umt,mplugowl}, InternVideo2 still grapples with limitations stemming from fixed input resolutions, sampling rates, and highly compressed tokens, which restrict its ability to express rich video information and capture fine-grained details.

The progressive learning scheme adopted by {\modelname} strikes a balance between model capabilities and training compute. While jointly learning the three optimization objectives simultaneously is computationally feasible, scalability becomes an issue when confronted with limited resources.

Although {\modelname} has demonstrated leading performance in some video understanding and reasoning benchmarks, it cannot guarantee an implicit world model that ensures consistency in visual reasoning. The inherent constraints imposed by fixed input representations, coupled with the complexity of visual reasoning tasks, present challenges in achieving a comprehensive and consistent understanding of the visual world.

\paragraph{Potential Biases.} We investigate the potential biases here. We focus on age, gender, and race distributions, as these are commonly recognized areas where bias can occur. 
We count keywords related to these categories in the used captions. Note that these synthetic captions may not fully reflect the truth of the corresponding videos, thereby creating a gap between our analysis and the actual reality. Here are the results of our analysis:
\begin{itemize}[leftmargin=*]
\item[$\bullet$] \textbf{Age}: The majority were about adults (86.99\%), followed by children (12.87\%) and barely any mentions of senior citizens (0.04\%).
\item[$\bullet$] \textbf{Gender}: 62.04\% pertained to men and 37.96\% pertained to women.
\item[$\bullet$] \textbf{Race}: 56.19\% are Asians, 23.04\% are Black people, 14.55\% are White people, 3.78\% are Middle Eastern people, and 2.43\% are Latin American people.
\end{itemize}

\section{Broader Impact}
It is important to acknowledge that, similar to other foundational models, {\modelname} has the potential to embed biases present in its training data and the associated models used during training, such as neural teachers \cite{videomae,chen2023internvl} and language models (LLMs) \cite{mistral,vicuna}. These biases may emerge due to a variety of factors, including the personal ideas, preferences, values, and perspectives of the data creators and the training corpus utilized.

The presence of biases in AI models can have societal implications and reinforce existing inequalities or prejudices. Biases within {\modelname} could manifest in the form of unfair or discriminatory outputs, potentially perpetuating social biases or stereotypes present in the training data. Consequently, it is crucial to carefully consider the potential impact of deploying {\modelname} in real-world applications and take proactive measures to mitigate biases and ensure fairness.